%!TEX program = xelatex

\def\numeromodulo{1}

\documentclass[12pt]{exam}

\def\ano{2021}
\def\semestre{2}
\def\disciplina{\'Algebra 1}
\def\turma{C}

\usepackage{caption}
\usepackage{amssymb}
\usepackage{amsmath,amsfonts,amsthm,amstext}
\usepackage[brazil]{babel}
% \usepackage[latin1]{inputenc}
\usepackage{graphicx}
\graphicspath{{/home/jfreitas/GitHub_Repos/Algebra-1/Pictures/}{D:/OneDrive - unb.br/imagens-latex/}}
\usepackage{enumitem}
\usepackage{multicol}
\usepackage{answers}
\usepackage{tikz,ifthen}
\usetikzlibrary{lindenmayersystems}
\usetikzlibrary[shadings]
\Newassociation{solucao}{Solution}{ans}
\newtheorem{exercicio}{}

\setlength{\topmargin}{-1.0in}
\setlength{\oddsidemargin}{0in}
\setlength{\textheight}{10.1in}
\setlength{\textwidth}{6.5in}
\setlength{\baselineskip}{12mm}

\extraheadheight{0.7in}
\firstpageheadrule
\runningheadrule
\lhead{
        \begin{minipage}[c]{1.7cm}
        \includegraphics[width=1.7cm]{unb.pdf}
        \end{minipage}%
        \hspace{0pt}
        \begin{minipage}[c]{4in}
          {Universidade de Bras{\'\i}lia} --
          {Departamento de Matem{\'a}tica}
\end{minipage}
\vspace*{-0.8cm}
}
% \chead{Universidade de Bras{\'\i}lia - Departamento de Matem{\'a}tica}
% \rhead{}
% \vspace*{-2cm}

\extrafootheight{.5in}
\footrule
\lfoot{\disciplina\ - \semestre$^o$/\ano}
\cfoot{}
\rfoot{P\'agina \thepage\ de \numpages}

\newcounter{exercicios}
\renewcommand{\theexercicios}{\arabic{exercicios}}

\newenvironment{questao}[1]{
\refstepcounter{exercicios}
\ifx&#1&
\else
   \label{#1}
\fi
\noindent\textbf{Exerc{\'\i}cio {\theexercicios}:}
}

\newcommand{\resp}[1]{
\noindent{\bf Exerc{\'\i}cio #1: }}
\newcommand{\nsub}{\varsubsetneq}
\newcommand{\vaz}{\emptyset}
\newcommand{\im}{{\rm Im\,}}
\newcommand{\sub}{\subseteq}
\newcommand{\n}{\mathbb{N}}
\newcommand{\z}{\mathbb{Z}}
\newcommand{\rac}{\mathbb{Q}}
\newcommand{\real}{\mathbb{R}}
\newcommand{\complex}{\mathbb{C}}
\newcommand{\cp}[1]{\mathbb{#1}}
\newcommand{\ch}{\mbox{\textrm{car\,}}\nobreak}
\newcommand{\vesp}[1]{\vspace{ #1  cm}}
\newcommand{\compcent}[1]{\vcenter{\hbox{$#1\circ$}}}
\newcommand{\comp}{\mathbin{\mathchoice
{\compcent\scriptstyle}{\compcent\scriptstyle}
{\compcent\scriptscriptstyle}{\compcent\scriptscriptstyle}}}

\begin{document}
    \begin{center}
    {\Large\bf \disciplina\ - Turma \turma\ -- \semestre$^{o}$/\ano} \\ \vspace{9pt} {\large\bf
        Lista de Exerc{\'\i}cios -- Módulo \numeromodulo}\\ \vspace{9pt} Prof. Jos{\'e} Ant{\^o}nio O. Freitas
    \end{center}
    \hrule

    \vspace{.6cm}

    \begin{center}
        \textit{Observa\c{c}\~oes: \begin{enumerate} \item A nota\c{c}\~ao $A \nsub B$ significa que $A \subset B$ e $A \ne B$, isto é, existe $x \in B$ com $x \notin A$. \item Seja $A$ um conjunto com uma quantidade finita de elementos, denotamos por $|A|$ a quantidade de elementos no conjunto $A$.\end{enumerate}}
    \end{center}
    \questao{} Dados os conjuntos $A = \{0,1,2\}$, $B = \{0,2,3\}$ e $C = \{0,1,2,3,4\}$, classifique as afirma\c{c}\~oes a seguir em verdadeira ou falsa, justificando:
    \begin{multicols}{2}
        \begin{enumerate}[label={\arabic*})]
            \item $A \sub B$

            \item $A \in C$

            \item $B \sub C$

            \item $A \in \{A, B, \{A, C\}\}$

            \item $\{A\} \sub \{B, \{A, B\}\}$

            \item $\{0,2\} \sub B$

            \item $\{0,2\} \sub C$

            \item $1 \sub C$

            \item $\{1,4\} \in C$

            \item $\{0,3\} \sub B$

            \item $B \sub A$

            \item $C \sub \{0, 1, 2, 3, 4, C, \{A, B, C\}\}$

            \item $\{\vaz\} \sub A$

            \item $\vaz \in C$

            \item $A \sub C$

            \item $x \neq \{x\}$

            \item $\emptyset = \{\emptyset\}$

            \item $\emptyset \in \{\emptyset\}$

            \item $\emptyset \sub \{A, B, \{C\}\}$

            \item $\{\} = \{0\}$

            \item $A \sub \emptyset$

            \item $\{\emptyset\} \sub \emptyset$

            \item $\{1, 3, 3\} = \{1, \{2, 3\}, 3\}$

            \item $\{a, b, c\} = \{A, B, C\}$

            \item $\{\{a\}, \{b\}\} = \{\{a\}, b\}$

            \item $\{1, 2, 3, 4, 5\} = \{4, 1, 3, 5, 2, 4, 5\}$

            \item $\{\} \sub \{\emptyset\}$

            \item $|\emptyset| = 1$

            \item $|\{x,x\}| = 2$

            \item $|A \cap \emptyset | = 0$

            \item $|\{A, B, C\}| = 11$

            \item $|\{A, B, C\}| = 3$

            \item $|\{A, \{B, C\}, C, \{A, B, C\}\}| = 4$

        \end{enumerate}
    \end{multicols}

    \questao{} D\^e exemplos de conjuntos n\~ao vazios $A$, $B$ e $C$ tais que:
    \begin{enumerate}[label={\alph*})]
        \item $A \sub B$, $C \sub B$ e $A \cap C = \emptyset$.

        \item $A \sub B$, $C \nsub B$ e $A \cap C = \emptyset$.

        \item $A \sub C$, $A \ne C$ e $B \cap C = \emptyset$.

        \item $A \sub (B \cap C)$, $B \sub C$, $B \ne C$ e $A \ne C$.

        \item $A \in B$, $B \sub C$ e $A \nsubseteq C$.

        \item $A \subsetneq B$, $B \in C$ e $A \in C$.

        \item $A \cap B \sub C$, $A \nsubseteq C$ e $B \nsubseteq C$.

        \item $A \cap C = \emptyset$, $A \sub B$ e $|B \cap C| = 3$.
    \end{enumerate}

    \vspace{.3cm}

    \questao{} Determine $A \cup B$ e $A \cap B$ se $A = \{x \in \real \mid x^2 = 1\}$ e $B = \z$.

    \vspace{.3cm}

    \questao{} Em cada um dos seguintes itens, determine se a afirma\c{c}\~ao \'e
    verdadeira ou falsa. Se for verdadeira, demonstre-a. Se for falsa, apresente um exemplo mostrando que a afirma\c{c}\~ao \'e falsa.
    \begin{enumerate}[label={\alph*})]
        \item Se $x \in A$ e $A \sub B$, ent\~ao $x \in B$.

        \item Se $A \nsub B$ e $B \sub C$, ent\~ao $A \nsub C$.

        \item Se $A \nsub B$ e $B \nsub C$, ent\~ao $A \nsub C$.

        \item Se $x \in A$ e $A \nsub B$, ent\~ao $x \notin B$.

        \item Se $x \in B$ e $A \nsub B$, ent\~ao $x \notin A$.

        \item Se $A \sub B$ e $x \notin B$, ent\~ao $x \notin A$.

        \item Se $A$, $B$ e $C$ s\~ao conjuntos, ent\~ao $A \cup (B \cap C) = (A \cup B) \cap C$.

        \item Se $A$, $B$ e $C$ s\~ao conjuntos, ent\~ao $A \cap (B \cup C) = (A \cup B) \cap C$.

        \item $A \cap B = B \cap A$.

        \item Se $A \cap B = A \cap C$, então $B = C$.

        \item Se $A \cup B = A \cup C$, então $B = C$.

        \item $A \cup B = B \cup A$.

        \item Se $A \cap B = A$, então $A \subset B$.

        \item Se $A \cup B = B$, então $A \subset B$.

        \item $A \cap B = A$ se, e somente se, $A \cup B = B$

        \item Se $A \sub B$, $B \sub C$ e $C \sub A$, então $A = B = C$.
    \end{enumerate}

    \vspace{.3cm}

    \questao{} Considere os seguintes conjuntos:
    \begin{align*}
        A &= \{x \in \z \mid x = 2m,\ m \in \z\}\\
        B &= \{y \in \z \mid y = 2(n - 1),\ n \in \z\}.
    \end{align*}
    Os conjuntos $A$ e $B$ são iguais? Justifique sua resposta.

    \vspace{.3cm}

    \questao{} Sejam $a$, $b \in \real$ com $a < b$. O intervalo fechado $[a, b]$ é o conjunto $\{x \in \real \mid a \le x \le b\}$. De modo semelhantes, o intervalo aberto $(a, b)$ é o conjunto $\{y \in \real \mid a < y < b\}$. Sejam $P = [3, 7]$, $Q = [7,9]$ e $R = [-3, 8]$. Descreva os seguintes conjuntos:
    \begin{enumerate}[label={\alph*})]
        \item $P \cap Q$.

        \item $P \cup Q$.

        \item $(P \cup Q) \cap R$.

        \item $P \cup (Q \cap R)$.
    \end{enumerate}

    \vspace{.3cm}

    \questao{} Sejam $A$ e $B$ conjuntos e suponha que $A$ não é um subconjunto de $B$. Quais das seguintes afirmações são verdadeiras? Encontre todas as respostas corretas.

    \begin{enumerate}[label={\roman*})]
        \item Se $x \in A$, então $x \notin B$.

        \item Se $x \in B$, ent~ao $x \in A$.

        \item Existe um elemento $x \in A$ tal que $x \notin B$.

        \item Existe um elemento $x \in B$, tal que $x \notin B$.
    \end{enumerate}
    \questao{} Demonstre que:
    \begin{enumerate}[label={\alph*})]
        \item Se $A \sub B$ e $C \sub D$, ent\~ao $A \cap C \sub B \cap D$.

        \item $A \cup B = \vaz$ se, e somente se, $A = \vaz$ e $B = \vaz$.

        \item $A \cup B = A \cap B$ se, e somente se,  $A = B$.

        \item $A \sub B$ se, e somente se,  $A \cap B = A$.

        \item $A \cup (B \cap (A \cup C)) = A \cup (B \cap C)$.

        \item $A \cup B = B$ se, e somente se, $A \subset B$.

        \item $(A \cup B) \cap C \sub A \cup (B \cap C)$.

        \item $A \sub B \cap C$ se, e somente se, $A \sub B$ e $A \sub C$.

        \item Se $A \sub B \cup C$ e $A \cap B = \emptyset$, então $A \sub C$.
    \end{enumerate}

    \questao{} Dê exemplos de conjuntos $A$, $B$ e $C$ tais que as interseções $A \cap B$, $A \cap C$ e $B \cap C$ não são o conjunto vazio mas $A \cap B \cap C$ é o conjunto vazio.

    \vspace{.3cm}

    \questao{} Determinar os elementos dos conjuntos $A$, $B$ e $E$ tais que:
    \[
        A \cap B = \{b, c\}, \quad C_E(A) = \{d, e, f\} \quad \mbox{e}\quad C_E(B) = \{a          , e, f\}.
    \]

    \vspace{.3cm}

    \questao{} Dados os conjuntos $A = \{a, b, c, d\}$, $B = \{c, d, e, f\}$, $C = \{d,       e, f, g\}$, determine:
    \begin{enumerate}[label={\alph*})]
        \item $A - (B - C)$ e $(A - B) - C$

        \item $(A \cap B) - (B \cup C)$ e $(A \cup C) - (A \cup B)$

        \item $A \cap (B - C)$ e $(A \cap B) - C$

        \item $A - (B \cup C)$ e $(A - B) \cap (B - C)$
    \end{enumerate}

    \questao{} Dados os conjuntos $E = \{1,2,3,4,5,a,b,c,d\}$, $A = \{1,2,4,d\}$ e $B =        \{a,2,4,b,5\}$, determine:
    \begin{multicols}{2}
        \begin{enumerate}[label={\alph*})]
            \item $A \cup B$

            \item $A \cup C_E(B)$

            \item $C_E(A) \cup B$

            \item $C_E(A \cup B)$
        \end{enumerate}
    \end{multicols}

    \vspace{.3cm}

    \questao{} Dados os conjuntos $A = \{3,6,9,12,15,18\}$, $B = \{2,3,5,7,11,13,17,19\}      $ e $C = \{1,2,3,4,5,\dots,20\}$ calcule:
    \begin{multicols}{3}
        \begin{enumerate}[label={\alph*})]
            \item $A - B$

            \item $A - C$

            \item $B - C$

            \item $(A - B) - C$

            \item $A - (B - C)$

            \item $(A \cup B) - C$

            \item $A - (B \cap C)$

            \item $(A \cup C) - B$

            \item $(B \cap C) - A$
        \end{enumerate}
    \end{multicols}

    \vspace{.3cm}

    \questao{} Determine os conjuntos $E$, $A$ e $B$ se:
    \begin{align*}
        C_E(A) &= \{2, 5, 9, 13, 18, 20\}\\
        C_E(B) &= \{2, 6, 18, 20\}\\
        A \cup B &= \{1, 5, 6, 9, 13, 14\}.
    \end{align*}

    \vspace{.3cm}

    \questao{} Demonstre que:
    \begin{enumerate}[label={\alph*})]

        \item Se $A \sub B$ e $C = B - A$, ent\~ao $A = B - C$.

        \item Se $A \cap B = \vaz$ e $A \cup B = C$, ent\~ao $A = C - B$.

        \item Se $A\cup B = E$, ent\~ao $C_E(A) \sub B$. \textit{(Aqui suponha que $A$,            $B \sub E$.)}

        \item Se $A \cap B = \vaz$, ent\~ao $A \cup C_E(B) = C_E(B)$. \textit{(Aqui suponha que $A$, $B \sub E$.)}

        \item Suponha que $A \subseteq E$. Mostre que $C_E((C_E(A) \cap E)) = A$.

        \item $A = B$ se, e somente se, $A - B = B - A$.

        \item $A \sub B$ se, e somente se, $A - B = \vaz$.

        \item $C_E(A) \sub C_E(B)$ se, e somente se, $A \cup B = A$. \textit{(Aqui suponha que $A$, $B \sub E$.)}

        \item $C_E(A) \sub C_E(B)$ se, e somente se, $A \cap B = B$. \textit{(Aqui suponha que $A$, $B \sub E$.)}

    \end{enumerate}

    \vspace{.3cm}

    \questao{} Demonstre as seguintes igualdades.
    \begin{enumerate}[label={\alph*})]
        \item $A \cup (C_E(A) \cap B) = A \cup B$. \textit{(Aqui suponha que $A$, $B \sub E$.)}

        \item $A \cap (C_E(A) \cup B) = A \cap B$. \textit{(Aqui suponha que $A$, $B \sub E$.)}

        \item $(A \cap C_E(B)) \cup (C_E(A) \cap B) = (A \cup B) \cap C_E(A \cap B)$. \textit{(Aqui suponha que $A$, $B \sub E$.)}
        \item $(A - B) - C = A - (B \cup C)$.

        \item $A \cup (B - C) = (A \cup B) - (C - A)$.

        \item $A \cap (B - C) = (A \cap B) - (A \cap C)$.

        \item $A - (B \cup C) = (A - B) \cap (A - C)$.

        \item $A - (B \cap C) = (A - B) \cup (A - C)$.

        \item $(A \cup B) - C = (A - C) \cup (B - C)$.

        \item $(A \cap B) - C = (A - C) \cap (B - C)$.

        \item $(A \cap B) \cap (A - B) = (A - B) \cap (B - A) = \vaz$.

        \item $(A - B) \cup (B - A) = (A \cup B) - (A \cap B)$.

        \item $A - (A - B) = A \cap B$.

        \item $(A - B) - (B - C) = A - B$.

        \item $(A - B) - B = A - B$.

        \item $A \cup (B \cap (A \cup C)) = A \cup (B \cap C)$.

        \item $(A \cap B) \cup C_E(A \cup C_E(B)) = B$. \textit{(Aqui suponha que $A$, $B \sub E$.)}

        \item $C_E( A \cap [B \cup C]) = C_E(A) \cup [C_E(B)\cap C_E(C)]$. \textit{(Aqui suponha que $A$, $B$, $C \sub E$.)}

        \item $C_E( A \cap B \cap C) = C_E(A) \cup C_E(B) \cup C_E(C)$. \textit{(Aqui suponha que $A$, $B$, $C \sub E$.)}

        \item $C_E( A \cup B \cup C) = C_E(A) \cap C_E(B) \cap C_E(C)$. \textit{(Aqui suponha que $A$, $B$, $C \sub E$.)}
    \end{enumerate}

    \vspace{.3cm}

    \questao{} Dados os conjuntos $A = \{1, 2, 3, r, s\}$, $B = \{r, s, 2, 3, 4\}$, $C = \{s, t, 4, 5\}$, $D = \{s, t, 5, 6\}$ e $E = \{1, 2, 3, 4, 5, 6, r, s, t\}$, determine:
    \begin{multicols}{2}
        \begin{enumerate}[label={\alph*})]
            \item $(A \times A) \cap (B \times B)$

            \item $A^2 \times C^2$

            \item $(A - B) \times (C - D)$

            \item $(A \cap B) \times (C \cap B)$

            \item $(A \cup B) \times (B \cup D)$

            \item $(A \times B) - (C \times D)$

            \item $(A - B) \times (C \cap D)$

            \item $(A - C) \times (B - D)$

            \item $(A - (C - D)) \times ((D - B) \cup A)$

            \item $C_E(A) \times C_E(B)$

            \item $C_{E\times E}(A \times B)$

            \item $(A - B) \times (C_E(B) - C_E(A))$
        \end{enumerate}

    \end{multicols}

    \vspace{.3cm}

    \questao{} Sejam $A$ e $B$ conjuntos tais que $A$, $B \subseteq E$. Mostre que $C_E(A) \times C_E(B) \subseteq C_{E \times E}(A \times B)$.

    \vspace{.3cm}

    \questao{} Seja $A$ e $B$ conjuntos tais que $A$, $B \subseteq E$. É verdade que
    \[
        C_{E\times E}(A \times B) = C_E(A) \times C_E(B)?
    \]
    Prove ou apresente um contra-exemplo.

    \vspace{.3cm}

    \questao{} Sejam $A$, $B$, $C$ e $D$ conjuntos quaisquer. A igualdade
    \[
        (A \cup B) \times (C \cup D) = (A \times C) \cup (B \times D)
    \]
    é válida? Justifique.

    \vspace{.3cm}

    \questao{} Sejam $A$, $B$, $C$ e $D$ conjuntos. Prove que:
    \begin{enumerate}[label={\alph*})]
        \item Se $A \sub C$ e $B \sub D$, ent\~ao $A \times B \sub C \times D$.

        \item $A \times (B \cup C) = (A \times B) \cup (A \times C)$

        \item $A \times (B \cap C) = (A \times B) \cap (A \times C)$

        \item $A \times B = \emptyset$ se, e somente se, $A = \emptyset$ ou $B = \emptyset$.

        \item $A \times B \ne \emptyset$ se, e somente se, $A \ne \emptyset$ e $B \ne \emptyset$.

        \item $(A \cap B) \times (C \cap D) = (A \times C) \cap (B \times D)$

        \item $(A \times B) - (C_E(A) \times C_E(B)) = [(A - C_E(A)) \times (B - C_E(B))] \cup [(A - C_E(A)) \times B]$  \textit{(Aqui suponha que $A$, $B \sub E$.)}

        \item $A \times (B - C) = (A \times B) - (A \times C)$

        \item $(A \times B) - (C \times D) = [A \times (B - D)] \cup [(A - C) \times B]$
    \end{enumerate}

    \vspace{.3cm}

    \questao{} Sejam $A$, $B$, $C$ e $D$ conjuntos quaisquer. É válido que
    \[
        (A \times B) - (C \times D) = (A - C) \times (B - D)?
    \]
    Prove ou dê um contra-exemplo.

    \vspace{.3cm}

    \questao{} Sejam $A$, $B$, $C$ e $D$ conjuntos.
    \begin{enumerate}[label={\alph*})]
        \item Mostre que $A$ e $B$ s\~ao disjuntos se, e somente se, para todo conjunto n\~ao vazio $C$, $A \times C$ e $B \times C$ s\~ao disjuntos.

        \item Suponha $A \ne \emptyset$ e $C \ne \emptyset$, com $A \ne C$. Mostre que $A \sub B$ e $C \sub D$ se, e somente se, $A \times C \sub B \times D$.
    \end{enumerate}

    \vspace{.3cm}

    \questao{} Sejam $X_1$, $X_1$, $Y_1$ e $Y_2$ subconjuntos contidos num conjunto $E$. Suponha que $X_1 \cup X_2 = E$, $Y_1 \cap Y_2 = \emptyset$, $X_1 \sub Y_1$ e $X_2 \sub Y_2$. Prove que $X_1 = Y_1$ e $X_2 = Y_2$.

    \vspace{.3cm}

    \questao{} Sejam $A$ e $B$ conjuntos quaisquer. Mostre que $A \subset B$ se, e somente se, $\mathcal{P}(A) \subset \mathcal{P}(B)$.

    \vspace{.3cm}

    \questao{} Sejam $A$ e $B$ conjuntos quaisquer. Mostre que $\mathcal{P}(A \cap B) = \mathcal{P}(A) \cap \mathcal{P}(B)$.

    \vspace{.3cm}

    \questao{} Sejam $A$ e $B$ conjuntos quaisquer. É verdade que
    \[
       \mathcal{P}(A \cup B) = \mathcal{P}(A) \cup \mathcal{P}(B)?
    \]
    Prove ou apresente um contra-exemplo.

    \vspace{.3cm}

    \questao{} Quais das rela{\c c}{\~o}es abaixo s{\~a}o rela{\c c}{\~o}es de equival{\^e}ncia sobre $E = \{a,b,c\}$?
    \begin{enumerate}[label={\alph*})]
        \item $R_1 = \{(a,a);(a,b);(b,a);(b,b);(c,c)\}$

        \item $R_2 = \{(a,a);(a,b);(b,a);(b,b);(b,c)\}$

        \item $R_3 = \{(a,a);(b,b);(b,c);(c,b);(a,c);(c,a)\}$

        \item $R_4 = E \times E$

        \item $R_5 = \vaz$
    \end{enumerate}

    \vspace{.3cm}

    \questao{} Determinar todas as rela{\c c}{\~o}es de equival{\^e}ncia
    $R$ sobre $A$ e os respectivos conjuntos quocientes $A/R$ para:
    \begin{enumerate}[label={\alph*})]
        \item $A=\{a\}$;

        \item $A=\{a,b\}$;

        \item $A=\{a,b,c\}$;

        \item $A=\{a,b,c,d\}$.
    \end{enumerate}

    \vspace{.3cm}

    \questao{} Seja $m \in \z$, $m > 1$. Defina $R \sub \z\times\z$ como
    \[
      R = \{(x,y) \in \z\times\z \mid x - y = km, \mbox{ para algum } k \in \z\}.
    \]
    Mostre que $R$ \'e uma rela\c{c}\~ao de equival\^encia sobre $\z$.

    \vspace{.3cm}

    \questao{} Para $x$, $y \in \real$ defina a relação $R$ como
    \[
        R = \{(x,y) \in \real \times \real \mid x - y = 2n\pi, \mbox{ para algum } n \in \z\}.
    \]
    Mostre que $R$ é uma relação de equivalência sobre $\real$.

    \vspace{.3cm}

    \questao{} Quais das seguintes senten{\c c}as definem uma rela{\c c}{\~a}o de equival{\^e}ncia no conjunto $A$ dado?
    \begin{enumerate}[label={\alph*})]
        \item $aRb$ se, e s{\'o} se, existe $k \in \n$ tal que $a - b = 3k$, $A = \n$.

        \item $aRb$ quando existe $k \in \n$ tal que $b = k a$, $A = \n$.

        \item $aRb$ quando $a \le b$, $A = \z$.

        \item $xRy$ quando $xy > 0$, $ A = \real$.

        \item $xRy$ quando $xy \ge 0$, $ A = \real$.

        \item $x \sim y$ quando $x + y$ \'e par, onde $x$, $y \in \z$.

        \item $x \sim y$ quando $x + y$ \'e {\'\i}mpar, onde $x$, $y \in \z$.

        \item $x \sim y$ quando $\dfrac{x}{y} \in \rac$, onde $x$, $y \in \real^* = \real - \{0\}$.

        \item $x \sim y$ quando $\dfrac{x}{y} \in \z$, onde $x$, $y \in \real^* = \real - \{0\}$.
    \end{enumerate}

    \vspace{.3cm}

    \questao{} Seja $\mathcal{A} = \{1,2,3,4\}$ e suponha que $R$ \'e uma rela\c{c}\~ao de equival\^encia em $\mathcal{A}$. Sabendo que $1R2$ e $2R3$ mostre que existem exatamente duas possibilidades para a rela\c{c}\~ao $R$, e descreva as duas.

    \vspace{.3cm}

    \questao{} Sejam $A$ e $B$ conjuntos. Definia a rela{\c c}\~ao $\mathcal{R}$ por
    \[
        A \mathcal{R} B \mbox{ quando } A \subseteq B.
    \]
    Ent\~ao $\mathcal{R}$ \'e uma rela{\c c}\~ao de equival\^encia?

    \vspace{.3cm}

    \questao{} Seja $A = \n \times \n^*$. Considere a seguinte
    rela{\c c}{\~a}o sobre $A$:
    \[
        (a,b) R (c,d) \mbox{ quando } a + b = c + d.
    \]
    Mostre que $R$ {\'e} uma rela{\c c}{\~a}o de equival{\^e}ncia sobre $A$.

    \vspace{.3cm}

    \questao{} Seja $A = \real^2$ e considere o conjunto definido por
    \[
      (a,b)R(c,d) \mbox{ quando } 2a - b = 2c - d.
    \]
    Mostre que $R$ \'e uma rela\c{c}\~ao de equival\^encia sobre $\real^2$.

    \vspace{.3cm}

    \questao{} Seja $A = \real^3$. Dados $(x, y, z)$, $(\alpha, \beta, \gamma) \in \real^3$, defina $(x, y, z) R (\alpha, \beta, \gamma)$ quando $z = \gamma$. Mostre que $R$ \'e uma rela\c{c}\~ao de equival\^encia sobre $\real^3.$

    \vspace{.3cm}

    \questao{} Seja $A = \real^3$. Dados $u = (x_1, y_1, z_1)$, $v = (x_2, y_2, z_2) \in \real^3$ defina
    \[
        u\cdot v = x_1x_2 + y_1y_2 + z_1z_2.
    \]
    Tome um elemento fixo $w = (\alpha, \beta, \gamma) \in \real^3$ e defina
    \[
        u \sim v \mbox{ quando } u \cdot w = v \cdot w.
    \]

    \begin{enumerate}[label={\alph*})]
        \item Mostre que $\sim$ \'e uma rela\c{c}\~ao de equival\^encia sobre $\real^3$.

        \item Encontre a classe de equivalência dos elementos $(1, 0, 0)$, $(0, 1, 0)$, $(0, 0, 1)$ e $(1, 1, 1)$.
    \end{enumerate}
    \vspace{.3cm}

    \questao{} Em $\z$ defina a relação $R$ por
    \[
        a R b \mbox{ se, e somente se, } |a + b| \mbox{ é par}
    \]
    para $a$, $b \in \z$.
    \begin{enumerate}[label={\alph*})]
        \item Mostre que $R$ é uma relação de equivalência em $\z$.

        \item Determine todas as classes de equivalência dessa relação.

        \item A qual classe de equivalência sua matrícula pertence.
    \end{enumerate}

    \vspace{.3cm}

    \questao{} Para pontos $(a, b)$, $(c, d) \in \real^2$ defina $(a, b) S (c, d)$ quando $a^2 + b^2 = c^2 + d^2$.
    \begin{enumerate}[label={\alph*})]
      \item Prove que $S$ \'e uma rela\c{c}\~ao de equival\^encia em $\real^2$.

      \item Liste todos os elementos no conjunto $\{(x, y) \in \real \mid (x, y) S (0, 0)\}$.

      \item Liste cinco elementos distintos no conjunto $\{(x, y) \in \real \mid (x, y) S (1, 0)\}$.
    \end{enumerate}

    \vspace{.3cm}

    \questao{} Sejam $E = \{-5, -4, -3, -2, -1, 0, 1, 2, 3, 5, 5\}$ e $R = \{(x, y) \in E \times E : x + |x| = y + |y|\}$.
    \begin{enumerate}[label={\alph*})]
        \item Mostre que $R$ \'e uma rela\c{c}\~ao de equival\^encia.

        \item Descreva o conjunto quociente $E/R$.
    \end{enumerate}
    \vspace{.3cm}

    \vspace{.3cm}

    \questao{} Considere a seguinte rela{\c c}{\~a}o sobre $\mathbb{C}$:
    \[
        (x+yi) R (r+si) \mbox{ quando } x^2+y^2=r^2+s^2.
    \]
    \begin{enumerate}[label={\alph*})]
        \item Mostre que $R$ {\'e} rela{\c c}{\~a}o de equival{\^e}ncia.

        \item Descreva a classe de equival{\^e}ncia de $2 - 3i$.
    \end{enumerate}

    \vspace{.3cm}

    \questao{} Considere a seguinte rela\c{c}\~ao sobre $\z$:
    \[
        x \sim y \mbox{ quando }  x^2 - y^2 = 4k, \mbox{ para algum } k \in \z.
    \]

    \begin{enumerate}[label={\alph*})]
        \item Mostre que $\sim$ {\'e} rela{\c c}{\~a}o de equival{\^e}ncia.

        \item Determine as classes de equivalência dessa relação.

        \item A qual classe de equivalência sua matrícula pertence.
    \end{enumerate}

    \vspace{.3cm}

    \questao{} Considera a rela\c{c}\~ao em $\real$ dada por
    \[
        a \sim b \mbox{ quando } a^2 + a = b^2 + b.
    \]
    \begin{enumerate}[label={\alph*})]
        \item Mostre que $\sim$ \'e uma rela\c{c}\~ao de equival\^encia em $\real$.

        \item Determine as classes de equival\^encia dessa rela\c{c}\~ao.
    \end{enumerate}

    \vspace{.3cm}

    \questao{} Seja $U$ um conjunto não vazio e considere $\mathcal{P}(U)$ o conjunto das partes de $U$. Fixado um conjunto $T$ defina a relação $\sim$ em $\mathcal{P}(U)$ por
    \[
        X \sim Y \mbox{ quando } X \cap T = Y \cap T,
    \]
    para $X$, $Y \in \mathcal{P}(U)$. Mostre que $\sim$ é uma relação de equivalência. Em particular, tome $U = \{1,2,3,4,5\}$ e $T = \{1,3\}$.
    \begin{enumerate}[label={\alph*})]
        \item Verdadeiro ou falso: $\{1,2,4\} \sim \{1,4,5\}$? Justifique.

        \item Verdadeiro ou falso: $\{1,2,4\} \sim \{1,3,4\}$? Justifique.

        \item Verdadeiro ou falso: $\{1,3,4\} \sim \{1,3,5\}$? Justifique.

        \item Encontre $\overline{\{1,5\}}$.

        \item Encontre $\overline{\{2,3\}}$.

    \end{enumerate}

    \vspace{.3cm}

    \questao{} Seja $\sim$ uma rela{\c c}{\~a}o sobre $\rac$ definida da seguinte forma:
    \[
        x \sim y \mbox{ quando } x - y \in \mathbb{Z}.
    \]
    \begin{enumerate}[label={\alph*})]
        \item Prove que $\sim$ {\'e} uma rela{\c c}{\~a}o de equival{\^e}ncia sobre $\rac$.

        \item Descreva a classe $\bar{1}$.

        \item Descreva a classe $\overline{1/2}$.
    \end{enumerate}

    \vspace{.3cm}

    \questao{} Defina
    \[
        H = \{3^m \mid m \in \z\} \mbox{ e } \rac^+ = \{x \in \rac \mid x > 0\}.
    \]
    Seja $R$ dado por
    \[
        R = \left\{(x,y) \in \rac^+\times \rac^+ : \frac{x}{y} \in H\right\}.
    \]
    \begin{enumerate}[label={\alph*})]
        \item Mostre que $R$ \'e uma rela\c{c}\~ao de equival\^encia em $\rac^+$.

        \item Determine a classe de equival\^encia de $5$.
    \end{enumerate}

    \vspace{.3cm}

    \questao{} Mostre que as seguintes condições definem uma relação de equivalência em $\z$ e encontre suas classes de equivalência:
    \begin{enumerate}[label={\alph*})]
        \item $mRn$ quando $3 | (m + 2n)$

        \item $mSn$ quando $5 | (2m + 3n)$.
    \end{enumerate}

    \vspace{.3cm}

    \questao{} A divisibilidade (ou seja, a rela{\c c}{\~a}o definida por $xRy$ se, e s{\'o}
    se, $x \mid y$) {\'e} uma rela{\c c}{\~a}o de equival{\^e}ncia sobre $\z$?

    \vspace{.3cm}

    \questao{} Seja $R$ a seguinte rela{\c c}{\~a}o sobre $\z^*$:
    \[
        x R y \mbox{ quando existem }  k, l \in \z \mbox{ tais que } y = kx \mbox{ e } x = ly.
    \]
    Mostre que $R$ {\'e} uma rela{\c c}{\~a}o de equival{\^e}ncia sobre $\z^*$.

    \vspace{.3cm}

    \questao{} Seja $R = \{ (x, y) \in \real^2 \mid x - y \in \rac\}$. Prove que $R$ {\'e} uma rela{\c c}{\~a}o de equival{\^e}ncia e descreva as classes representadas por $1/2$ e $\sqrt{2}$.

    \vspace{.3cm}

    \questao{} Seja $A$ um conjunto n\~ao vazio. Suponha que $R \sub A \times A$ \'e tal que:
    \begin{enumerate}[label={\alph*})]
        \item Para todo $x \in A$, $(x,x) \in R$;

        \item Para todos $x$, $y$, $z \in A$ se $(x, y) \in R$ e $(x,z) \in R$, ent\~ao $(y,z) \in R$.
    \end{enumerate}
    Mostre que $R$ \'e uma rela\c{c}\~ao de equival\^encia sobre $A$.

    \vspace{.3cm}

    \questao{} Sejam $R$ e $S$ relações de equivalência em um conjunto $A$. Defina $T = R \cap S$. Assim $aTb$ se, e somente se, $aRb$ e $aSb$. Mostre que $T$ é uma relação de equivalência em $A$.

    \vspace{.3cm}

    \questao{} É verdade que a união de duas relações de equivalência ainda é uma relação de equivalência?
\end{document}
